% !PRINTED 144
\weberpage{144}
%<<CONT
\weq{10}{\frac{1}{\mathrm{dn}(u+v)} = \frac{\mathrm{dn}\,u\,\mathrm{dn}\,v + \varkappa^2\mathrm{sn}\,u\,\mathrm{sn}\,v\,\mathrm{cn}\,u\,\mathrm{cn}\,v}{\mathrm{dn}^2u\,\mathrm{dn}^2v + \varkappa^2\varkappa'^2\mathrm{sn}^2u\,\mathrm{sn}^2v},}
\weq{11}{\frac{\mathrm{dn}(u+v)}{\mathrm{cn}(u+v)} = \frac{\mathrm{dn}\,u\,\mathrm{dn}\,v\,\mathrm{cn}\,u\,\mathrm{cn}\,v + \varkappa'^2\mathrm{sn}\,u\,\mathrm{sn}\,v}{\mathrm{cn}^2u\,\mathrm{cn}^2v - \varkappa'^2\mathrm{sn}^2u\,\mathrm{sn}^2v},}
\weq{12}{\frac{\mathrm{cn}(u+v)}{\mathrm{dn}(u+v)} = \frac{\mathrm{dn}\,u\,\mathrm{dn}\,v\,\mathrm{cn}\,u\,\mathrm{cn}\,v - \varkappa'^2\mathrm{sn}\,u\,\mathrm{sn}\,v}{\mathrm{dn}^2u\,\mathrm{dn}^2v + \varkappa^2\varkappa'^2\mathrm{sn}^2u\,\mathrm{sn}^2v}.}
Man kann noch mannigfache andere Formeln auf dieselbe Weise herleiten, unter denen wir die folgenden drei anf\"uhren, die sich durch Division von \S\,22 (4), (3), (1) durch (2) ergeben.
\weq{13}{\mathrm{sn}(u+v)\,\mathrm{sn}(u-v) = \frac{\mathrm{sn}^2u - \mathrm{sn}^2v}{1 - \varkappa^2\mathrm{ns}^2u\,\mathrm{sn}^2v},}
% [sic] Printed ns in the first factor; the parallel denominators (14), (15) print sn.
\weq{14}{\mathrm{cn}(u+v)\,\mathrm{cn}(u-v) = \frac{\mathrm{cn}^2u\,\mathrm{cn}^2v - \varkappa'^2\mathrm{sn}^2u\,\mathrm{sn}^2v}{1 - \varkappa^2\mathrm{sn}^2u\,\mathrm{sn}^2v},}
\weq{15}{\mathrm{dn}(u+v)\,\mathrm{dn}(u-v) = \frac{\mathrm{dn}^2u\,\mathrm{dn}^2v + \varkappa^2\varkappa'^2\mathrm{sn}^2u\,\mathrm{sn}^2v}{1 - \varkappa^2\mathrm{sn}^2u\,\mathrm{sn}^2v}.}
Die wichtigsten unter diesen Formeln sind (5), (7), (9), aus denen sich, wenn auch durch weitl\"aufige Rechnungen, die \"ubrigen alle herleiten lassen. Der \"Ubersicht halber setzen wir sie noch einmal her:
\weq{16}{\begin{aligned}
\mathrm{sn}(u+v) &= \frac{\mathrm{sn}\,u\,\mathrm{cn}\,v\,\mathrm{dn}\,v + \mathrm{sn}\,v\,\mathrm{cn}\,u\,\mathrm{dn}\,u}{1 - \varkappa^2\mathrm{sn}^2u\,\mathrm{sn}^2v},\\
\mathrm{cn}(u+v) &= \frac{\mathrm{cn}\,u\,\mathrm{cn}\,v - \mathrm{sn}\,u\,\mathrm{sn}\,v\,\mathrm{dn}\,u\,\mathrm{dn}\,v}{1 - \varkappa^2\mathrm{sn}^2u\,\mathrm{sn}^2v},\\
\mathrm{dn}(u+v) &= \frac{\mathrm{dn}\,u\,\mathrm{dn}\,v - \varkappa^2\mathrm{sn}\,u\,\mathrm{sn}\,v\,\mathrm{cn}\,u\,\mathrm{cn}\,v}{1 - \varkappa^2\mathrm{sn}^2u\,\mathrm{sn}^2v}.
\end{aligned}}
Indem man je zwei dieser Gleichungen kombiniert, kann man ihnen unter anderen die Formen geben:
\weq{17}{\begin{aligned}
\mathrm{cn}(u+v)\,\mathrm{dn}\,u\,\mathrm{dn}\,v - \mathrm{dn}(u+v)\,\mathrm{cn}\,u\,\mathrm{cn}\,v &= -\varkappa'^2\mathrm{sn}\,u\,\mathrm{sn}\,v\\
\mathrm{dn}(u+v)\,\mathrm{sn}\,u\,\mathrm{cn}\,v - \mathrm{sn}(u+v)\,\mathrm{dn}\,u &= -\mathrm{cn}\,u\,\mathrm{sn}\,v,\\
\mathrm{sn}(u+v)\,\mathrm{cn}\,u - \mathrm{cn}(u+v)\,\mathrm{sn}\,u\,\mathrm{dn}\,v &= \mathrm{dn}\,u\,\mathrm{sn}\,v,\\
\mathrm{cn}(u+v)\,\mathrm{cn}\,u + \mathrm{sn}(u+v)\,\mathrm{sn}\,u\,\mathrm{dn}\,v &= \mathrm{cn}\,v.
\end{aligned}}

Wir wenden die Formeln (16) zun\"achst an zur Feststellung der Periodizit\"at der elliptischen Funktionen, die sich nat\"urlich auch aus \S\,21 herleiten l\"a\ss t. Setzt man in (16) $v = \pm K$, $v = K + iK'$, so folgt aus \S\,42 (16):
%CONT>>
